% Same body/content as the IEEE submission; only the template wrapper differs.
\documentclass[11pt]{article}
\usepackage[margin=1in]{geometry}
\usepackage{cite}
\usepackage{amsmath,amssymb,amsfonts}
\usepackage{cancel}
\usepackage{algorithmic}
\usepackage{graphicx}
\usepackage{xcolor}
% Review annotations (author review only): \rev{...} prints a red bracketed note.
\newcommand{\rev}[1]{\textcolor{red}{[REVIEW: #1]}}
\usepackage{textcomp}
\usepackage[ruled,linesnumbered]{algorithm2e}
\usepackage{caption}
\usepackage{subcaption}
\usepackage{float}
\usepackage{siunitx}
\usepackage{booktabs}
\usepackage{wasysym}
\usepackage{url}
\usepackage{makecell}
\usepackage{booktabs}
\usepackage{threeparttable}
\usepackage{multirow}
\usepackage{appendix}
\usepackage[verbose]{placeins}


% Acronyms
\usepackage[acronym]{glossaries}
\newacronym{sdf}{SDF}{Signed Distance Function}
\newacronym{aabb}{AABB}{Axis Aligned Bounding Box}
\newacronym{mc}{MC}{Marching Cubes}
\newacronym{mil}{MIL}{Mean Intercept Length}
\newacronym{da}{DA}{Degree of Anisotropy}
\newacronym{smi}{SMI}{Structural Model Index}
\newacronym{ccc}{CCC}{Concordance Correlation Coefficient}
\newacronym{tbth}{Tb.Th}{Trabecular Thickness}
\newacronym{tbsp}{Tb.Sp}{Trabecular Separation}
\newacronym{connd}{Conn.D.}{Connectivity Density}
\newacronym{tbn}{Tb.N.}{Trabecular Number}
\newacronym{nrmse}{nRMSE}{Normalized Root Mean Squared Error}
\newacronym{roi}{ROI}{Region of Interest}
\newacronym{gui}{GUI}{Graphical User Interface}
\newacronym{pla}{PLA}{Polylactic Acid}
\newacronym{knn}{KNN}{K-nearest Neighbors}
\newacronym{microct}{micro-CT}{Micro-computed Tomography}
\newacronym{oat}{OAT}{One-at-a-time}

% bibtex
\def\BibTeX{{\rm B\kern-.05em{\sc i\kern-.025em b}\kern-.08em
    T\kern-.1667em\lower.7ex\hbox{E}\kern-.125emX}}
\begin{document}

\title{ScaffoldBuilder: A free software tool for creating bone scaffolds. \\ User's Guide}
\author{Konstantinos Risvas}

\maketitle

\section{Introduction}
This document is provided as a user guide for the software. We present the main concepts and capabilities as well as a step-by-step tutorial, that users can follow to familiarize themselves with the tool. 

\section{Main Concepts}

The tool can be used to create 3D scaffold models for porous / foam structures, such as the trabecular bone. The generation process is dictated by a set of parameters to adjust pore separation, thickness, plate / rod balance and anisotropy. The generated mesh is a manifold watertight mesh that can be either exported as a surface mesh or as a binary image.

\begin{enumerate}

\item Decide container domain between the following
\begin{enumerate}
\item Box
\item Cylinder
\item Random Convex Mesh
\end{enumerate}

\item Generate points inside the container:
\begin{enumerate}
\item Random seeding: Generate randomly the points inside the container. Determine the number of points using the GUI.
\item Uniform Poisson sampling: Generate points inside the container by applying a minimum distance constraint between them.
\item Varied Poisson sampling: Generate points inside the container by applying a distance constraint between them. The constraint is applied as a distance function either from the container surface, a custom plane, or a unique point.
\end{enumerate}

\item Decide the parameters of the scalar field for surface extraction with \gls{mc}:
\begin{enumerate}
\item Voxel Size: The voxel size decides the resolution of the generation grid. It is crucial to capture thin structure features according to the Nyquist limit $f_N = f_s / 2$, where $f_s$ is the sampling frequency and $f_N$ the maximum value of that feature that can be measured without distortion. For example, for an iso-level value of 0.15 a voxel size of at least 0.075 should be used with values in the range $0.05 - 0.075$ more suitable for a fine mesh generation.
\item Scaffold Thickness: This the iso-level value that the algorithm of \gls{mc} will extract the surface. It is the thickness of the scaffold struts and walls. Users can apply either a uniform thickness value or one generated by a distance function from a reference (point, plane container's walls) and with a custom modulation (linear, quadratic, sigmoid).
\item Anisotropy: Users can specify global and local anisotropy sources. Each one of those is equipped with a covariance matrix $\mathbf{C}_{m}=\mathbf{R}_{m}^{\!\top}\mathbf{S}_{m}^{2}\mathbf{R}_{m}$. $\mathbf{R}_{m}$ is the rotation matrix that can be used to apply abstract directions on respective material region and $\mathbf{S}_{m}$ is the stretch diagonal matrix $\mathrm{diag}(s_x,s_y,s_z)$ with the $s_k$ values being the principal stretches. Each source has its own area of influence that is blended with a global anisotropy background with $C_{bg}$.
\item Smoothness: Users can adjust the smoothness levels of the geometry in a variety of ways. Initially, a Taubin smoothing of the scalar field on the voxel level can be used to smooth out "jigsaw" patterns and other aliasing artifacts. Moreover, we provide a smoothing of the sharp Voronoi polyhedra and the vertex junctions with a polynomial minimum smoothing.   
\end{enumerate}

\item Calibration Mode: Users can supply target values for the medical descriptors of \gls{tbth}, porosity, and \gls{da}. Then a nested loop performs successive 1D root finding for each metric using a secant method that tries to adjust the iso-level, openness and principal stretch values until the targets are satisfied. Users can decide which of the three variables to optimize.

\item Create Scaffold: Based on the applied settings the scaffold is generated and rendered. 

\item Export Scaffold: The generated scaffold can be exported as a surface mesh (.stl, .vtk formats) or as a 3d volume (stack of 2d slices - useful for using with ImageJ).

\item I/O utilities: Apart from surface or binary formats scaffolds can be extracted/loaded with custom formats.
\begin{itemize}
    \item .scaf: This can be used to extract a single scaffold entity. It stores the entire scalar field, the list of parameters and clinical metrics as well as the anisotropy sources and the applied container and seed generator.
    \item .sbproj: This format is used to save/load an entire project. This means that each scaffold is stored as a .scaf file and the main .sbproj file contains all generated projects (containers, generators, anisotropy sources, \glspl{roi}).
\end{itemize}

\end{enumerate}

\section{Installation}

A binary is provided in the GitHub page: \textcolor{red}{add link when ready}. The user can extract the contents in the desired folder and execute the provided application file (ScaffoldBuilder.exe).

\section{GUI Presentation}

In Fig.~\ref{fig:gui_overview}, an overview of the panel is presented. At the top, is the main menu bar that provides access to I/O utilities and rendering and algorithm settings. The toolbar provides a row of buttons that can be used to control the distinct steps of the generation process, from container selection to geometry creation as well as estimation of clinical measurements. On the left, there is the control panel that provides access to the properties of each generated entity, such as the scaffolds, containers, seed generators, anisotropy sources, and \glspl{roi}. Specific tools, such as the "Plane Cut", are rendered in individual windows. At the bottom, there is a console that provides users with information about each action and generation step. In the middle of the screen, the main rendering window is used to render generated objects that have a visual model. Each one of those has customizable appearance in terms of color and opacity. All windows can be moved and docked, thus, the entire \gls{gui} is completely reconfigurable.

\begin{figure}[ht!]
    \centering
    \includegraphics[width=1.0\linewidth]{Figures/gui_overview.png}
    \caption{Overview of the Graphical User Interface (GUI). It consists of a main menu bar, a toolbar for specific actions, a rendering window for Visualization, a control panel for the generated objects, and a console window for conveying messages to the user. }
    \label{fig:gui_overview}
\end{figure}

\subsection{Main Menu Bar}

The main menu bar consists of six sub--menus. The first is the "File" that contains the I/O utilities for project and scaffold loading/exportation. The "View" menu contains rendering options for visualization objects, such as normals, mesh edges, bounding box etc. The "Window" menu can be accessed to hide/show the main windows of the \gls{gui}. In the "Settings" context menu, users can adjust algorithm and rendering settings. The "Add" menu can be used to create anisotropy sources and \glspl{roi}. Finally, the "Tools" bar consists of the "Cut With Plane" and "Cut With ROI" tools.

\begin{figure}[h!]
    \centering
    \includegraphics[width=0.75\linewidth]{Figures/gui_main_menu_bar.png}
    \caption{The main menu bar provides a control panel for configuring visualization settings, performing I/O utilities, and accessing specific tools.}
    \label{fig:gui_menu_bar}
\end{figure}
\FloatBarrier

\subsection{Toolbar}

The "Toolbar" is a row of buttons that control the following:

\renewcommand{\labelenumii}{\arabic{enumi}.\arabic{enumii}}
\renewcommand{\labelenumiii}{\arabic{enumi}.\arabic{enumii}.\arabic{enumiii}}
\renewcommand{\labelenumiv}{\arabic{enumi}.\arabic{enumii}.\arabic{enumiii}.\arabic{enumiv}}
\begin{enumerate}
    \item Container generation
    \begin{enumerate}
        \item Box Container
        \item Cylindrical Container
        \item Abstract Container
    \end{enumerate}
    \item Seed generation
    \begin{enumerate}
        \item Random Seed Generator
        \item Uniform Seed Generator
        \item Varied Seed Generator
    \end{enumerate}
    \item Clinical descriptors
    \begin{enumerate}
        \item Estimate Tb.Th.
        \item Estimate Tb.Sp.
        \item Estimate tortuosity
        \item Estimate \gls{da}
        \item Estimate \gls{smi}
        \item Estimate \gls{connd}
        \item Estimate \gls{tbn}
    \end{enumerate}
    \item Mesh Generation
    \begin{enumerate}
        \item Mesh Creation
        \item Mesh Update
    \end{enumerate}
    \item Mesh Processing
    \begin{enumerate}
        \item Translation
        \item Scaling
        \item Taubin Smoothing
        \item Surface Simplification
    \end{enumerate}
    \item Undo / Redo actions
    \begin{enumerate}
        \item Undo action
        \item Redo action
    \end{enumerate}
\end{enumerate}

\begin{figure}[h!]
    \centering
    \includegraphics[width=1.0\linewidth]{Figures/gui_toolbar.png}
    \caption{The main menu bar provides a control panel for configuring visualization settings, performing I/O utilities, and accessing specific tools.}
    \label{fig:gui_toolbar}
\end{figure}
\FloatBarrier

\subsection{Object Window}

\subsection{Properties Window}

\subsection{Visualizer Window}

\section{Tutorial 1 - Box Scaffold Creation}
In this tutorial, we present the required steps for creating a scaffold inside a box container. 

\subsection{Container Creation}

Start by pressing the textured button on the left side of the toolbar.

\begin{figure}[h!]
    \centering
    \includegraphics[width=1.0\linewidth]{Figures/select_container.png}
    \caption{Step1. Press the box container button. This will cause a popup window to open.}
    \label{fig:box_button}
\end{figure}

Pressing the button will open the Box Container popup window. Leave default values for now. This will create a box of $10 \times 10 \times 10$~mm that is centered in the box center $\mathbf{O} = [5, 5, 5]$. Edit the "Name" field and type "box container". This will change the name of the object in the control panel. Press the "Create" button.  

\begin{figure}[h!]
    \centering
    \includegraphics[width=0.5\linewidth]{Figures/box_container_settings.png}
    \caption{Step 2. Box Container Settings. Adjust object name, side size and center of box. Press "Create" to generate the container.}
    \label{fig:box_container_settings}
\end{figure}
\FloatBarrier

This will create the container object that will be added to the "Containers" list of the control panel "Objects". Initially, the container is set as hidden. It can be visualized by pressing right-click on the name and then "Show". The "Properties" window is automatically updated to show the information and properties of each active object in the "Objects" window.

\begin{figure}[h!]
    \centering
    \includegraphics[width=0.5\linewidth]{Figures/box_container_properties.png}
    \caption{Step 3. Visualize container model and inspect settings on the "Properties" window. Users can update the visualization properties and the actual size and position of the container. Pressing "Update" will regenerate the container object.}
    \label{fig:box_container_properties}
\end{figure}
\FloatBarrier

\subsection{Seed Generation}

Subsequently, from the toolbar locate and press the "Random Seed Generator" button.

\begin{figure}[h!]
    \centering
    \includegraphics[width=0.5\linewidth]{Figures/random_seed_generator_button.png}
    \caption{Step 4. Press the suggested button to create a Random Seed Generator.}
    \label{fig:random_seed_generator_button}
\end{figure}
\FloatBarrier

A popup will open where the list of available containers is displayed and the number of seeds can be specified. 

\begin{enumerate}
    \item Select the generated "box container" from the list
    \item Edit the name to "random generator"
    \item Set the number of seeds to 500
    \item Press Generate
\end{enumerate}

\begin{figure}[h!]
    \centering
    \includegraphics[width=0.5\linewidth]{Figures/random_seed_generator_settings.png}
    \caption{Step 5. Specify the generator settings. }
    \label{fig:random_seed_generator_settings}
\end{figure}
\FloatBarrier

The object will be added to the "Generators" list under the "Objects" window. Users can regenerate the seeds by adjusting the desired properties. The seeds are rendered as a point cloud in the "Visualizer" window.

\begin{figure}[h!]
    \centering
    \includegraphics[width=1.0\linewidth]{Figures/random_seed_generator_model.png}
    \caption{Generator Properties and Model}
    \label{fig:random_seed_generator_model}
\end{figure}
\FloatBarrier

\subsection{Scaffold Creation}

The next step is to create the actual scaffold. Press the button that is render with a scaffold texture.

\begin{figure}[h!]
    \centering
    \includegraphics[width=1.0\linewidth]{Figures/scaffold_button.png}
    \caption{Step 6. Press the "Scaffold Generation" button. This will open the "Scaffold Creator" panel.}
    \label{fig:scaffold_button}
\end{figure}
\FloatBarrier

The panel provides access to all specific parameters for generating the scalar field and these are organized in thematic tabs. The main tab is used to set the model's name, select the proper container and generator, define the desired voxel size for the generation grid. Also, users can set the measurement voxel size, that will be used to resample the grid and estimate the clinical descriptors that use a voxel--based calculation.

\begin{enumerate}
    \item Set the scaffold name to "Box Scaffold"
    \item Leave the voxel size and measurement voxel size to the default values
    \item Select the box container and the random generator from the lists
    \item Uncheck "Calibrate Porosity"
\end{enumerate}

\begin{figure}[h!]
    \centering
    \includegraphics[width=1.0\linewidth]{Figures/scaffold_generator_main.png}
    \caption{Step 7. In the main tab edit the scaffold's name, the generation voxel size and the measurement voxel size. Select the container and the generator, and if desired, set a calibration target for porosity.}
    \label{fig:scaffold_main}
\end{figure}
\FloatBarrier

Next,
\begin{enumerate}
    \item Click the "Thickness" tab.
    \item Uncheck "Calibrate Thickness"
    \item Leave thickness mode to uniform and iso--level to the default value
    \item Set openness to $0.75$. This will lead to more open faces and a geometry with higher - porosity.
    \item Leave spread to 0.5 to have a balance between rod--plate structures
\end{enumerate}

\begin{figure}[h!]
    \centering
    \includegraphics[width=1.0\linewidth]{Figures/scaffold_generator_thickness.png}
    \caption{Step 8. In the "Thickness" tab users can select the thickness mode (uniform, varied) the iso--level value, and the target \gls{tbth} for calibration. Also, the openness value that determines the degree of fenestration and the spread value that sets the rod--plate heterogeneity can be defined in this window.}
    \label{fig:scaffold_thickness}
\end{figure}
\FloatBarrier

Subsequently, let's apply a global anisotropy to the scaffold. Do the following:
\begin{enumerate}
    \item Select the "Anisotropy" tab
    \item Uncheck the "Calibrate Anisotropy" button for now
    \item Set the stretch along x to $1.5$
    \item Set the direction vector to $[0, 1, 1]$ and hit "Normalize". This will normalize the direction to $[0, 0.707, 0.707]$. The x stretch is the principal stretch along the direction specified by this normal vector.
\end{enumerate}

\begin{figure}[h!]
    \centering
    \includegraphics[width=1.0\linewidth]{Figures/scaffold_generator_anisotropy.png}
    \caption{Step 8. In the "Thickness" tab define the anisotropic settings. The stretch x is the principal stretch value along the direction vector specified in the "Direction" field.}
    \label{fig:scaffold_anisotropy}
\end{figure}
\FloatBarrier

Finally, let's define some smoothing parameters:

\begin{enumerate}
    \item Select the "Smoothness" tab
    \item Set the Taubin iterations to 100.
    \item Check the "Round Cell Edges" options and set the $k_e$ value to 0.2
    \item Add the container's edges to the solid phase. The value is automatically set to the selected iso--level.
\end{enumerate}

\begin{figure}[h!]
    \centering
    \includegraphics[width=1.0\linewidth]{Figures/scaffold_generator_smoothness.png}
    \caption{Step 8. In the "Smoothness" tab define the smoothness settings. Here, users can adjust the Taubin filtering parameters, as well as the minimum polynomial smoothing for junctions and edges. Additionally, we provide options for including the container's edges and the scaffold's boundaries to the solid phase.}
    \label{fig:scaffold_smoothness}
\end{figure}
\FloatBarrier

We are ready to create our first scaffold. Hit "Generate". A yellow-filling bar informs about the completion stage of the generation process. Upon creation, the scaffold object is added to the "Scaffolds list". The "Properties" panel provides access to the active scaffold's properties and settings. Also, a field with the clinical descriptors is provided.

\begin{figure}[h!]
    \centering
    \includegraphics[width=1.0\linewidth]{Figures/scaffold_properties_panel.png}
    \caption{Step 8. When a scaffold object is selected from the "Objects" window then its properties are rendered under the "Properties" panel. Users can update the generation settings and also overview the estimated clinical descriptors. When a metric is rendered with an orange color, it means that its version is outdated and needs re-estimation. Some of the metrics are automatically calculated during the generation process.}
    \label{fig:scaffold_properties_panel}
\end{figure}
\FloatBarrier

The model is rendered in the "Visualizer" window. The appearance can be adjusted from the properties panel. Change the color and see the results. We can also render the edges and normals of the surface mesh.

\begin{enumerate}
    \item Change the color from the "Appearance" field
    \item In the main menu bar click on "View" and select "Show Scaffold Face Normals" and "Show Scaffold Edges". These action can be also performed using the keys "N" and "M", respectively. 
\end{enumerate}

\begin{figure}[h!]
    \centering
    \includegraphics[width=0.5\linewidth]{Figures/change_view.png}
    \caption{Step 8. The "View" panel of the main menu bar can be used to change rendering settings.}
    \label{fig:change_view}
\end{figure}
\FloatBarrier

\begin{figure}[h!]
    \centering
    \includegraphics[width=0.5\linewidth]{Figures/change_color.png}
    \caption{Step 8. Change the object's color from the "Appearance" field.}
    \label{fig:change_color}
\end{figure}
\FloatBarrier

\begin{figure}[h!]
    \centering
    \includegraphics[width=0.5\linewidth]{Figures/normals_edges.png}
    \caption{Step 8. Users can check normals and edges of the scaffold by enabling the respective rendering settings. Here the normals are displayed as green thin lines and the edges as black. Observe, also that the container's edges are included the solid geometry, like a cage.}
    \label{fig:normals_edges}
\end{figure}
\FloatBarrier

\subsection{Metric Estimation}

Let's now estimate the \gls{tbth} of the scaffold. 
\begin{enumerate}
    \item On the toolbar identify and press the "Estimate DA" button. This will fire a popup with the settings
    \item Leave the voxel size at its default value. When the window opens the value of this field is set to "Measurement Voxel Size" we already have defined during the scaffold creation process.
    \item Leave the "Analyze Entire Scaffold" option checked. This will use the entire scalar field for estimating \gls{da}.
    \item Hit "Estimate"
    \item Observe the metric in the "Properties" panel. It is now white, denoting that it matches the current scaffold version
\end{enumerate}

\begin{figure}[h!]
    \centering
    \includegraphics[width=0.5\linewidth]{Figures/estimate_da.png}
    \caption{Step 8. By pressing the "Estimate DA" button a popup opens that users can set the algorithm settings for estimating \gls{da}.}
    \label{fig:estimate_da}
\end{figure}
\FloatBarrier

Upon completion the metric field has an updated color:

\begin{figure}[h!]
    \centering
    \includegraphics[width=0.5\linewidth]{Figures/estimate_da_change.png}
    \caption{Step 8. When \gls{da} is re-estimated the properties are updated. The white color informs that the \gls{da} version matches the current scaffold version.}
    \label{fig:estimate_da_change}
\end{figure}
\FloatBarrier

\subsection{Update Scaffold}

We can update the morphology of the generated scaffold by either updating the container or/and the generator settings or by adjusting the scaffold's parameters from the "Properties" panel.

Let's do both. 

\begin{enumerate}
    \item Select the "Box Container" from the "Containers" list.
    \item Set the size along the z axis to $5$~mm and adjust the z--coordinate of the center to be at half the size along the z axis.
    \item Press "Update".
    \item Select the "Random Generator" from the "Seed Generators" list.
    \item Press "Update".
    \item Select again the "Box Scaffold" from the "Scaffolds" list.
    \item In the "Thickness" tab set thickness to $0.2$~mm.
    \item Set openness to $1.0$.
    \item Set spread to $0.0$ this will disable the spread knob and leave plate--rod heterogeneity unchanged.
    \item From the toolbar locate the "Update" button which displays a green "Play" texture. Press it.
    \item The scaffold will be regenerated. Observe that the \gls{da} metric is now orange--ish. Thus, we have to re--estimate it to comply with the current mesh version. 
    \item Observe the material direction and how the model is mainly consisted of struts.
\end{enumerate}

\begin{figure}[h!]
    \centering
    \includegraphics[width=1.0\linewidth]{Figures/scaffold_update_gen_con.png}
    \caption{Step 8. Update the box container and the random seed generator.}
    \label{fig:scaffold_update_gen_con}
\end{figure}
\FloatBarrier

\begin{figure}[h!]
    \centering
    \includegraphics[width=0.8\linewidth]{Figures/scaffold_update_thickness.png}
    \caption{Step 8. Update the thickness properties under the "Thickness" tab and hit the "Update" button. This will regenerate the scaffold with the new container, generator, and thickness parameters.}
    \label{fig:scaffold_update_thickness}
\end{figure}
\FloatBarrier

\begin{figure}[h!]
    \centering
    \includegraphics[width=0.8\linewidth]{Figures/box_scaffold_model_updated.png}
    \caption{Step 8. The model is updated. Observe the new dimensions along z and the material direction along $[0, 0.707, 0.707]$. Also, the \gls{da}metric is now orange, denoting that it needs to be estimated again to match the current scffold version.}
    \label{fig:box_scaffold_model_updated}
\end{figure}
\FloatBarrier

\subsection{Exporting}

Scaffold models can be exported in three different ways. The first one is software--specific. Users can store all relevant parameters, metrics, and related entities of a scaffold in a ".scaf" format. Notice that this will save the currently selected scaffold from the "Objects" menu.

\begin{enumerate}
    \item From the main menu bar select "File" and "Export Scaffold".
    \item Create a new directory pressing the "+" button, type a new name and hit "OK".
    \item The new directory will be opened. Type a new name "box\_scaffold.scaf" and press "OK".
\end{enumerate}

\begin{figure}[h!]
    \centering
    \includegraphics[width=0.5\linewidth]{Figures/box_scaffold_export_scaf.png}
    \caption{test}
    \label{fig:box_scaffold_export_scaf}
\end{figure}
\FloatBarrier

\begin{figure}[h!]
    \centering
    \includegraphics[width=0.25\linewidth]{Figures/box_scaffold_export_scaf_2.png}
    \caption{test}
    \label{fig:box_scaffold_export_scaf_2}
\end{figure}
\FloatBarrier

\begin{figure}[h!]
    \centering
    \includegraphics[width=0.5\linewidth]{Figures/box_scaffold_export_scaf_3.png}
    \caption{test}
    \label{fig:box_scaffold_export_scaf_3}
\end{figure}
\FloatBarrier

There are two ways to save a scaffold as a surface mesh (".stl", ".vtk") or a binary image (".nrrd"). 

\begin{enumerate}
    \item Navigate to the main menu bar select "File" and "Save Scaffold as geometry" or "Save Scaffold as binary image". This will save the current active scaffold.
    \item Locate the name of the scaffold in the "Objects" list and right--click on it. There are two options: "Export as Mesh" and "Export as binary image". 
\end{enumerate}

\begin{figure}[h!]
    \centering
    \includegraphics[width=0.5\linewidth]{Figures/box_scaffold_export_geo_bin_1.png}
    \caption{test}
    \label{fig:box_scaffold_export_geo_bin_1}
\end{figure}
\FloatBarrier

\begin{figure}[h!]
    \centering
    \includegraphics[width=0.5\linewidth]{Figures/box_scaffold_export_geo_bin_2.png}
    \caption{test}
    \label{fig:box_scaffold_export_geo_bin_2}
\end{figure}
\FloatBarrier

Users can also export the generation parameters and estimated metrics in a ".csv" file using the right context menu of the scaffold object under the "Objects" list. Similar steps can be followed to load a scaffold from a ".scaf" file or from a geometry file (".stl", ".vtk"). The advantage of the first file format is that the scaffold is fully reconfigurable whereas the geometry model can only be loaded for visual inspection.

\section{Tutorial 2 - Calibration}

In this tutorial, we are going to load the previously created scaffold, reconfigure it and apply the three calibration steps for creating a geometry that satisfies targeted values for \gls{da}, \gls{tbth}, and porosity.

\begin{enumerate}
    \item From the main menu bar and the "View" tab click on "Reset Scene". This will clear all the objects from the scene and reset the current project. From the "File" tab click on "Load Scaffold", navigate to the folder where you exported the scaffold of Tutorial 1, select "box\_container.scaf" and click on "OK".
    \item The scaffold will be regenerated starting from the \gls{mc} step.
    \item Select the scaffold from the "Objects" list and select the "Main" tab. Check the "Calibrate Porosity" and set the target at $75\%$.
    \item Press the "Update" button. The scaffold will be regenerated, and the porosity target will be achieved.
\end{enumerate}

\begin{figure}[h!]
    \centering
    \includegraphics[width=1.0\linewidth]{Figures/tut_2_calibrate_porosity.png}
    \caption{test}
    \label{fig:tut_2_calibrate_porosity}
\end{figure}
\FloatBarrier

Let's also apply the calibration for \gls{tbth} and \gls{da}. 

\begin{enumerate}
    \item Click the "Thickness" tab.
    \item Activate the "Calibrate Thickness" checkbox. 
    \item Set the target \gls{tbth} to $0.25$~mm.
    \item Open the "Anisotropy" tab.
    \item Check "Calibrate Anisotropy".
    \item Set the target \gls{da} to 1.7
    \item Press "Update".
\end{enumerate}

\begin{figure}[h!]
    \centering
    \includegraphics[width=1.0\linewidth]{Figures/tut_2_calibrate_thickness.png}
    \caption{test}
    \label{fig:tut_2_calibrate_thickness}
\end{figure}
\FloatBarrier

\begin{figure}[h!]
    \centering
    \includegraphics[width=1.0\linewidth]{Figures/tut_2_calibrate_anisotropy.png}
    \caption{test}
    \label{fig:tut_2_calibrate_anisotropy}
\end{figure}
\FloatBarrier

The calibration process will take some time to complete. It required $\approx 278$~seconds for the machine that was used to create this tutorial. 

\begin{enumerate}
    \item Check the updated metrics. 
    \item Check the updated geometry.
\end{enumerate}

\begin{figure}[h!]
    \centering
    \includegraphics[width=1.0\linewidth]{Figures/tut_2_calibrate_check.png}
    \caption{test}
    \label{fig:tut_2_calibrate_check}
\end{figure}
\FloatBarrier

Save the scaffold as "box\_scaffold\_calibrated.scaf". You can also estimate the other clinical metrics.

\section{Tutorial 3 - Create Scaffold Inside Femur Container}

In this tutorial, we are going to create a porous geometry inside an abstract container. 

\subsection{Container}

\begin{enumerate}
    \item In the toolbar click on the "Abstract Container" button. Navigate to the tutorial folder and select the "femoral\_head\_container.stl". The container will be inserted in the "Containers" list.
    \item Select the container and adjust its scaling by 10 and press "Update".
    \item Observe the updated size of the container in the "Visualizer" window.
\end{enumerate}    

\begin{figure}[h!]
    \centering
    \includegraphics[width=1.0\linewidth]{Figures/tut_3_container.png}
    \caption{test}
    \label{fig:tut_3_container}
\end{figure}
\FloatBarrier

\subsection{Seed Generation}

Let's create a seed cloud using a varied 3D Poisson sampling. 

\begin{enumerate}
    \item Click in the varied Poisson generator.
    \item Select the container.
    \item Set the distance function as "Distance from Container".
    \item Define the modulation as quadratic.
    \item Set the starting radius to $0.6$ and the ending radius to $0.3$ and the transition distance to $1.0$.
    \item Press Create.
    \item Observe the generated points.
\end{enumerate}

\begin{figure}[h!]
    \centering
    \includegraphics[width=1.0\linewidth]{Figures/tut_3_generate_seeds.png}
    \caption{test}
    \label{fig:tut_3_generate_seeds}
\end{figure}
\FloatBarrier

\subsection{Locally Varying Anisotropy}

We are going to enforce a varying anisotropy to our geometry. 

\begin{enumerate}
    \item From the main menu bar select "Add" and click "Add Anisotropy Source"
    \item In the "Create Anisotropy Source" window change the position to $[0.5, 1.0, -2.0]$.
    \item Change sigma value to $1.0$.
    \item Change the principal stretch to $2.0$.
    \item Set the direction vector to $[0.0, 1.0, -1.0]$ and hit "Normalize".
    \item Press "Add".
\end{enumerate}

\begin{figure}[h!]
    \centering
    \includegraphics[width=1.0\linewidth]{Figures/tut_3_anisotropy.png}
    \caption{test}
    \label{fig:tut_3_anisotropy}
\end{figure}
\FloatBarrier


\subsection{Scaffold Generation}

Now it's time to generate the porous geometry. Click on the "Scaffold" icon and the generator will pop up. We are going to apply a varied thickness and the anisotropy blending.


\begin{enumerate}
    \item Press again the button with the geometry icon.
    \item In the "Scaffold Creator" window under the "Main" tab, in the "Name" field type "femur scaffold".
    \item Select the container "femoral\_head.stl" and the "varied generator".
    \item Uncheck "Calibrate Porosity"
\end{enumerate}

\begin{figure}[h!]
    \centering
    \includegraphics[width=1.0\linewidth]{Figures/tut_3_scaffold_1.png}
    \caption{test}
    \label{fig:tut_3_scaffold_1}
\end{figure}
\FloatBarrier

Subsequently, let's set the thickness parameters:

\begin{enumerate}
    \item Select the "Thickness" tab. 
    \item Uncheck the "Calibrate Thickness" option.
    \item Select the "Varied Thickness" option.
    \item Set "Start Thickness" to $1.0$ and "End Thickness" to $0.15$. Set the "Transition Distance" to $2.0$.
    \item Select "Distance From Container" as reference.
    \item Define the modulation as "Sigmoid".
    \item Set "Openness" to $0.8$.
\end{enumerate}

\begin{figure}[h!]
    \centering
    \includegraphics[width=1.0\linewidth]{Figures/tut_3_scaffold_thickness.png}
    \caption{test}
    \label{fig:tut_3_scaffold_thickness}
\end{figure}
\FloatBarrier

Next, we are going to define anisotropy.

\begin{enumerate}
    \item Select the "Anisotropy" tab. 
    \item Uncheck the "Calibrate Anisotropy" option.
    \item Set "Stretch Y" to $1.5$.
    \item Ensure that "source 1" is checked.
\end{enumerate}

\begin{figure}[h!]
    \centering
    \includegraphics[width=1.0\linewidth]{Figures/tut_3_scaffold_anisotropy.png}
    \caption{test}
    \label{fig:tut_3_scaffold_anisotropy}
\end{figure}
\FloatBarrier

Finally, define some smoothing parameters:

\begin{enumerate}
    \item Select the "Smoothness" tab. 
    \item Set Taubin smoothing iterations to $100$.
    \item Enable "Round Cell Edges".
    \item Set "k" to $0.15$.
    \item Hit "Generate".
\end{enumerate}

\begin{figure}[h!]
    \centering
    \includegraphics[width=1.0\linewidth]{Figures/tut_3_scaffold_smoothness.png}
    \caption{test}
    \label{fig:tut_3_scaffold_smoothness}
\end{figure}
\FloatBarrier

The scaffold will be generated and render in the "Visualizer" window.

\begin{enumerate}
    \item From the main menu bar select "Tools" and "Cut With Plane".
    \item Set the tool's normal to $[-1.0, 0.0, 0.0]$.
\end{enumerate}

Observe that around the femoral head the anisotropy is influenced by the added source whereas the rest of the mesh is influenced by the global anisotropy background with stretch $1.5$ along the y--axis.

\begin{figure}[h!]
    \centering
    \includegraphics[width=0.5\linewidth]{Figures/tut_3_scaffold_cut.png}
    \caption{test}
    \label{fig:tut_3_scaffold_cut}
\end{figure}
\FloatBarrier

\begin{figure}[h!]
    \centering
    \includegraphics[width=1.0\linewidth]{Figures/tut_3_scaffold_final.png}
    \caption{test}
    \label{fig:tut_3_scaffold_final}
\end{figure}
\FloatBarrier

The generated mesh has a total of $\approx 600000$ faces. The number of vertices and faces depends on the resolution of the grid. We have implemented a mesh simplification method using Quadric Error Metrics~\cite{garland1997surface}.

\begin{enumerate}
    \item Press the "Surface Simplification" button.
    \item In the "Simplify Mesh" window set "Target Triangles" to $0$.
    \item Leave the target percentage to $25\%$.
    \item Press "Generate".
    \item Press the "M" from keyboard.
    \item Observe that now the mesh consists of larger triangles and that the number of faces is $148112$. If you replicate the experiments the numbers may slightly differ.
\end{enumerate}

\begin{figure}[h!]
    \centering
    \includegraphics[width=1.0\linewidth]{Figures/tut_3_scaffold_qem.png}
    \caption{test}
    \label{fig:tut_3_scaffold_qem}
\end{figure}
\FloatBarrier

\section{Conclusions}

In this document, we presented three tutorials that demonstrate the software's capabilities. 

\newpage
\begin{thebibliography}{00}

\bibitem{garland1997surface}
M.~Garland and P.~S. Heckbert, ``Surface simplification using quadric error metrics,'' in \emph{Proceedings of the 24th annual conference on Computer graphics and interactive techniques}, pp.~209--216, ACM Press/Addison-Wesley Publishing Co., 1997.

\end{thebibliography}

\end{document}